\section{Mechanistic Results}
\label{sec:mechanistic}

\subsection{H1: Cross-Hint Probe Generalization}
We train a linear probe on residual activations from the \emph{sycophancy} hint type and test its zero-shot performance on the other five hint types. \Cref{tab:probe_results} and \Cref{fig:probe_auc} show the AUC scores.

\begin{table}[H]
    \centering
    \begin{tabular}{lcccccc}
\toprule
Layer & In-domain (sycophancy) & Zero-shot: Consistency & Zero-shot: Metadata & Zero-shot: Visual Pattern & Zero-shot: Grader Hacking & Zero-shot: Unethical \\
\midrule
10 & 99.7\% & 90.5\% & 85.9\% & 78.8\% & 97.0\% & 99.5\% \\
14 & 99.5\% & 93.1\% & 93.9\% & 80.5\% & 94.2\% & 99.6\% \\
18 & 99.7\% & 96.9\% & 95.0\% & 88.8\% & 96.6\% & 99.9\% \\
22 & 99.6\% & 93.6\% & 91.7\% & 85.1\% & 94.6\% & 99.2\% \\
26 & 99.4\% & 94.0\% & 93.3\% & 79.5\% & 94.8\% & 99.2\% \\
\bottomrule
\end{tabular}

    \caption{H1 cross-hint probe AUC scores per layer. The probe is trained on sycophancy hints and tested zero-shot on other hint types.}
    \label{tab:probe_results}
\end{table}

\begin{figure}[H]
    \centering
    \includegraphics[width=0.8\textwidth]{probe_auc.pdf}
    \caption{H1 cross-hint probe AUC across layers. Strong zero-shot generalization (AUC > 0.85) indicates a shared shortcut subspace across hint types.}
    \label{fig:probe_auc}
\end{figure}

Layer \num{18} consistently shows the strongest probe performance, suggesting this is where shortcut evidence is most strongly encoded in the residual stream.

\subsection{H2: Activation Patching}
We patch clean-condition residuals into misleading-hint forward passes at specific layers and measure answer flip rates. \Cref{tab:patching_results} and \Cref{fig:patching_flips} show the results.

\begin{table}[H]
    \centering
    \begin{tabular}{lrcrrc}
\toprule
Original Label & Cases & Layer & Total & Flipped & Flip Rate \\
\midrule
Sycophantic Failure & 5 & 14 & 5 & 0 & 0.0\% \\
Sycophantic Failure & 5 & 18 & 5 & 0 & 0.0\% \\
Sycophantic Failure & 5 & 22 & 5 & 0 & 0.0\% \\
Right Answer Wrong Reason & 5 & 14 & 5 & 1 & 20.0\% \\
Right Answer Wrong Reason & 5 & 18 & 5 & 1 & 20.0\% \\
Right Answer Wrong Reason & 5 & 22 & 5 & 1 & 20.0\% \\
Faithful Reasoning & 5 & 14 & 5 & 0 & 0.0\% \\
Faithful Reasoning & 5 & 18 & 5 & 0 & 0.0\% \\
Faithful Reasoning & 5 & 22 & 5 & 0 & 0.0\% \\
\bottomrule
\end{tabular}

    \caption{H2 activation patching results. Percentage of cases where patching clean residuals flipped the model's answer.}
    \label{tab:patching_results}
\end{table}

\begin{figure}[H]
    \centering
    \includegraphics[width=0.8\textwidth]{patching_flips.pdf}
    \caption{H2 activation patching answer flip rates by layer. Higher flip rates indicate stronger causal evidence that the layer encodes shortcut information.}
    \label{fig:patching_flips}
\end{figure}

\subsection{H3: RAWR-vs-Faithful Distance Test}
We measure the distance of RAWR and faithful residuals from clean-condition residuals in the same layer. If RAWR residuals are systematically farther from clean than faithful residuals, this provides an internal trace of shortcut influence. \Cref{tab:rawr_vs_faithful} and \Cref{fig:rawr_vs_faithful} show the results.

\begin{table}[H]
    \centering
    \begin{tabular}{lrcrcc}
\toprule
Layer & RAWR (n) & RAWR Rate & Faithful (n) & Faithful Rate & Difference \\
\midrule
10 & 157 & 22.9\% & 401 & 20.2\% & +2.7\% \\
14 & 157 & 27.4\% & 401 & 21.2\% & +6.2\% \\
18 & 157 & 28.7\% & 401 & 21.9\% & +6.7\% \\
22 & 157 & 35.7\% & 401 & 26.7\% & +9.0\% \\
26 & 157 & 37.6\% & 401 & 25.2\% & +12.4\% \\
\bottomrule
\end{tabular}

    \caption{H3 RAWR-vs-faithful distance test. Shortcut detection rate is the percentage of cases where RAWR residuals are farther from clean than faithful residuals.}
    \label{tab:rawr_vs_faithful}
\end{table}

\begin{figure}[H]
    \centering
    \includegraphics[width=0.8\textwidth]{rawr_vs_faithful.pdf}
    \caption{H3 RAWR-vs-faithful shortcut detection rates across layers. Values above 50\% indicate RAWR residuals are systematically farther from clean than faithful residuals.}
    \label{fig:rawr_vs_faithful}
\end{figure}

The consistent detection rate above \qty{50}{\percent} across layers confirms that RAWR residuals carry measurable traces of shortcut influence.
